% ============================================================
\chapter{Modul 3: Schleifen}
% ============================================================

Ein Roboter-Kontrollprogramm läuft nicht einmal -- es läuft in einer Endlosschleife: immer wieder Sensoren lesen, Entscheidung treffen, Motoren steuern. Schleifen sind das Herzstück autonomer Systeme.

\section{Lektion 8: Warum Schleifen?}

\begin{analogie}
Eine Produktionsanlage prüft jede Sekunde: Temperatur im Soll? Druck im Soll? Fehler? Diese Überprüfung passiert nicht einmal -- sie passiert 86.400-mal am Tag. Ohne Schleifen müsste man dieselbe Prüfung 86.400-mal aufschreiben.
\end{analogie}

\begin{lstlisting}
# Ohne Schleife: unmoeglich fuer viele Wiederholungen
motor_schritt()
motor_schritt()
motor_schritt()
# ... 1000-mal tippen?

# Mit Schleife: elegant
for i in range(1000):
    motor_schritt()
\end{lstlisting}

\section{Lektion 9: Die \texttt{for}-Schleife}

\begin{lstlisting}
# Grundform: N-mal ausfuehren
for i in range(5):
    print(f"Schritt {i}")
# Ausgabe: Schritt 0, Schritt 1, ..., Schritt 4

# Mit Start und Stop
for i in range(10, 101, 10):   # 10, 20, 30, ..., 100
    print(f"Geschwindigkeit: {i}%")

# Ueber eine Liste
sensoren = ["IR-vorne", "IR-links", "IR-rechts", "Ultraschall"]
for sensor in sensoren:
    print(f"Pruefe {sensor}...")
    # hier wuerde man den echten Sensorwert lesen

# Mit enumerate: Index und Wert gleichzeitig
wegpunkte = [(0,0), (1,0), (1,1), (0,1)]
for i, (x, y) in enumerate(wegpunkte):
    print(f"Wegpunkt {i}: x={x}, y={y}")
\end{lstlisting}

\begin{merksatz}
\texttt{range(stop)}: 0 bis stop-1. \texttt{range(start, stop)}: start bis stop-1. \texttt{range(start, stop, step)}: mit Schrittweite. Die obere Grenze ist \textbf{nie} enthalten.
\end{merksatz}

\begin{praxis}
Wenn ein Roboter eine Trajektorie (eine Liste von Zielpunkten) abfährt, iteriert das Programm mit \texttt{for} über die Liste der Wegpunkte -- genau wie im letzten Beispiel.
\end{praxis}

\begin{uebung}[title=Übung 3.1: Sensor-Scan]
Schreibe eine \texttt{for}-Schleife, die einen simulierten 180-Grad-Scan macht: Winkel von 0 bis 180 in 15-Grad-Schritten. Für jeden Winkel: gib Winkel und einen simulierten Abstandswert aus (du kannst ihn berechnen: z.\,B. \texttt{abstand = 1.0 / (winkel + 1) * 180}).
\end{uebung}

\section{Lektion 10: Die \texttt{while}-Schleife}

\begin{lstlisting}
# Hauptkontrollschleife eines Roboters
akku = 100
position = 0.0

while akku > 10:
    # Sensor lesen (simuliert)
    abstand = 0.5   # in echtem Code: sensor.read()

    if abstand < 0.20:
        print("Hindernis! Stoppe.")
        break   # Schleife sofort verlassen

    # Fahren
    position += 0.1
    akku -= 2
    print(f"Position: {position:.1f} m | Akku: {akku}%")

print("Schleife beendet.")
if akku <= 10:
    print("Akku leer -- Ladestation anfahren.")
\end{lstlisting}

\begin{merksatz}
\texttt{while BEDINGUNG:} laeuft solange die Bedingung \texttt{True} ist. \texttt{break} beendet die Schleife sofort. \texttt{continue} ueberspringt den Rest des aktuellen Durchlaufs. Vergisst man, die Bedingung zu veraendern: Endlosschleife. Stopp: \texttt{Strg+C}.
\end{merksatz}

\begin{lstlisting}
# Die echte Roboter-Hauptschleife sieht so aus:
laufe = True

while laufe:
    sensor_daten = lese_sensoren()
    entscheidung = verarbeite(sensor_daten)
    fuehre_aus(entscheidung)

    if notaus_gedrueckt():
        laufe = False
\end{lstlisting}

\begin{ros}
In ROS2 heißt diese Endlosschleife \texttt{rclpy.spin(node)}. Sie läuft bis das Programm beendet wird und verarbeitet dabei alle eingehenden Messages. Das ist das direkte Äquivalent zur \texttt{while True:}-Hauptschleife.
\end{ros}

\begin{uebung}[title=Projektaufgabe Modul 3: Simulated Drive Loop]
Schreibe eine simulierte Hauptkontrollschleife für deinen Roboter:
\begin{itemize}
    \item \texttt{while}-Schleife läuft solange Akku > 15\% und kein Notaus
    \item Jede Iteration: Akku sinkt um 1, Position wächst um 0.2
    \item Zufälliger Abstand: nutze \texttt{import random; abstand = random.uniform(0.1, 2.0)}
    \item Bei Abstand < 0.25: Warnung ausgeben, Richtung ändern (Variable \texttt{richtung})
    \item Nach 50 Iterationen automatisch stoppen
    \item Abschlussbericht: gefahrene Strecke, Restakku
\end{itemize}
\end{uebung}
